\chapter{検証実験（RQ2・RQ3）}
本章では，OshiFitの有用性を検証するために行う6週間の検証実験について述べる．実験の目的と仮説，倫理的配慮，参加者，群の定義，手続き，測定項目，分析計画を述べた後，結果と考察を示す．\todo{実験実施後に5.8節・5.9節を記入する}

\section{目的と仮説}
本実験の目的は，RQ2「6週間の利用で，すでに良好な人（群A），プランBを選んだ人（群B），プランAを選んだ人（群C）の幸福度と推し疲れは，それぞれの群の中でどう変化し，群の間でどう異なるか」と，RQ3「利用者は自分の推し疲れの原因に気づき，特に行動を変えた人はシステムを有用と評価するか」に答えることである．次の仮説を立てる．

\begin{description}
\item[H1] 群Bでは，開始時に比べて6週目の推し疲れが低下し，幸福度（WHO-5）が維持または向上する．
\item[H2] 開始時の状態と性格を調整したうえで，群Bは群Cに比べて6週目の推し疲れが低く，幸福度が高い．
\item[H3] 群Bの中で，プランを実行した程度が高い人ほど，推し疲れの低下が大きい（用量反応関係）．
\item[H4] 群Bと群Cの利用者の多くが，OshiFitによって自分の推し疲れの原因に気づいたと回答し，特に群Bの利用者はOshiFitを有用と評価する．
\end{description}

群Aは開始時に良好な状態にある人であるため，6週間で大きな変化は想定しないが，平均への回帰によってわずかに低下する可能性がある．

\section{倫理的配慮}
本実験は，大阪工業大学の研究倫理審査の承認を得てから実施する．\todo{要記入：審査委員会の名称，承認番号，承認日}本研究は人を対象とする医学系研究には該当しないと指導教員と判断し，臨床試験の登録は行わない．

参加者には，研究の目的，内容，所要時間，データの扱い，参加の任意性と途中で取りやめられることを画面で説明し，同意を得てから開始する．氏名や連絡先などの個人を特定できる情報は収集しない．写真と顔の特徴点は端末の外に送らず，保存しない．データはFirebase（東京リージョン）に匿名の参加IDで保存し，閲覧は研究者（筆者と指導教員）に限る．参加の取りやめを選んだ参加者の記録は削除する．WHO-5や推し疲れの得点が一定の基準を下回った場合には相談先を表示する（4.11節）．実験の終了後には，全参加者に研究の目的と結果の概要を説明し（デブリーフィング），すべての機能を利用できるようにする．

\section{参加者}
推し活を日常的に行っている18歳以上の人を対象とし，大学の連絡網やSNSを通じて募集する．目標とする参加者数は約240名とした．これは，6週目までの離脱と群Aの割合を見込んだうえで，群Bと群Cの比較に必要な人数を確保するためである．\todo{要記入：実際の募集期間，募集方法，参加者数，属性，離脱者数}

\section{群の定義}
本実験では，参加者を無作為に群に割り付けることはしない．群は，開始時の状態と本人の選択によって決まる．
\begin{description}
\item[群A（すでに良好）] 開始時に6軸のうち5軸以上が60点以上，かつWHO-5の素点が13点以上の人．プランは提示せず，6週目に再度回答してもらう．
\item[群B（プランBを選択）] 群Aに当てはまらず，2つの未来を見比べたうえでプランBを選んだ人．
\item[群C（プランAを選択）] 群Aに当てはまらず，2つの未来を見比べたうえでプランA（今のまま）を選んだ人．
\end{description}

無作為化を行わない理由は，趣味である推し活に対して，研究の都合で変更を強制したり禁止したりすることが，設計指針D1（「今のまま」を常に有効な選択肢とする）に反するためである．その代わりに，群の差を解釈する際には，開始時の状態と性格による違いを統計的に調整し，結果は効果ではなく「関連」として解釈する．

\section{手続き}
\begin{enumerate}
\item \textbf{開始時（約10分）}：参加者は説明を読んで同意し，67項目の質問紙に回答し，アバターを作成する．OshiFitは6軸の得点とその算出根拠を表示し，群を判定する．
\item \textbf{選択}：群Aに当てはまらない参加者は，3つのプランとプランA・Bの見込みを見て，プランBを選ぶか，プランAを選ぶかを決める．
\item \textbf{6週間}：群Bと群Cの参加者は，週1回以上，振り返りに回答する（約2分）．群Aの参加者は，期間中の利用は不要とする．
\item \textbf{6週目}：全員がWHO-5を含む終了調査に回答する．群Bと群Cの参加者は，気づき，使いやすさ，有用性に関する項目にも回答する．
\item \textbf{終了後}：全員にデブリーフィングを行う．
\end{enumerate}

習慣の形成には個人差が大きく，数週間から数か月を要することが知られている\cite{lally2010}．本実験では，参加者の負担と修士研究の期間を考慮して6週間とし，習慣の定着ではなく，プランの実行と推し疲れ・幸福度の変化の関係を捉えることを目的とする．

\section{測定項目}
\subsection{主要な結果指標}
\begin{itemize}
\item \textbf{幸福度}：WHO-5\cite{topp2015,inagaki2013}の素点（0〜25）．開始時と6週目に測定する．
\item \textbf{推し疲れ（総合）}：Q44〜Q51の8項目から算出した0〜100点（表\ref{tab:axes}）．開始時と6週目に測定する．
\end{itemize}

\subsection{副次的な結果指標}
\begin{itemize}
\item 6軸それぞれの得点（開始時と6週目）
\item 週ごとの推し疲れ4項目と推し活中の幸福感
\item プランの実行率（行った$=1$，一部行った$=0.5$，行わなかった$=0$の6週間の平均）
\item 気づき（「OshiFitを使って，自分の推し疲れの原因に気づいた」など）
\item 使いやすさと有用性（群B・C）\todo{要確認：終了調査の項目の最終版（付録D）}
\end{itemize}

\section{分析計画}
分析の計画は，データの収集を始める前に定めた．無作為化を行わない設計における群の比較方法（共分散分析，傾向スコアによる重み付け，多重比較の補正）については，データサイエンス学科の濱田悦夫教授から統計解析に関する助言を得た．

\subsection{各群の前後比較}
各群について，開始時と6週目のWHO-5と推し疲れを，対応のある$t$検定とWilcoxonの符号順位検定で比較し，効果量$d_z$を求める．

\subsection{3群の比較}
6週目の値を目的変数とし，群，開始時の値，性格5因子を説明変数とする共分散分析（ANCOVA）を行う．群Aを基準として，群Bと群Cとの差を推定する．群Aは開始時の状態で選ばれているため，差は効果ではなく関連として解釈する．

\subsection{群Bと群Cの比較}
プランBを選ぶかどうかは本人の選択であるため，選択のしやすさに関わる要因によって群Bと群Cの間に違いが生じうる．そこで，性格，開始時の6軸・WHO-5・推し疲れ，推し活の頻度から，プランBを選ぶ確率（傾向スコア）\cite{rosenbaum1983,hoshino2009}を推定し，その逆数で重み付けした回帰によって，開始時の値を調整したうえで群Bと群Cの6週目の値を比較する．傾向スコアで調整しても測定されていない違いは残るため，結論は関連までとする．

\subsection{用量反応関係}
群Bの中で，プランの実行率と推し疲れ・幸福度の変化量の関係を回帰分析で調べる．実行率が高いほど推し疲れの低下が大きいという関係が見られれば，プランが変化に関わっていることの根拠が強まる．

\subsection{多重比較}
主要な比較（WHO-5と推し疲れ$\times$各群の前後比較と群Bと群Cの比較）には，Holm法による補正を行う．

\subsection{正解データによる想定改善量の更新}
実験の終了後に1回だけ，群Bの正解データ（プランを実行し，ストレスが減り，幸福感が下がっていない週）を用いてプランの想定改善量を更新する．プランごとに，正解データが1週以上ある参加者について狙った軸の実際の変化量（6週目$-$開始時）の平均を求め，
\begin{equation}
\text{新しい想定改善量}=\frac{10\times\text{これまでの想定改善量}+n\times\text{観測された平均}}{10+n}
\end{equation}
とする（$n$はそのプランの該当人数）．人数が少ないプランは元の値から大きく動かない．更新した値をモデルv1.1とし，研究終了後の利用に用いる．実験期間中はv1.0のまま変えない．

\section{結果}
\todo{実験実施後に記入}

\subsection{参加者の内訳}
\todo{要記入：募集人数，同意者数，6週目まで回答した人数，各群の人数と属性}

\begin{table}[t]
\centering
\jetcaption{各群の参加者の内訳と開始時の得点（実験後に記入）}{Participants and baseline scores by group (to be completed)}
\label{tab:groups}
\small
\begin{tabular}{@{}lccc@{}}
\toprule
 & 群A & 群B & 群C \\
\midrule
人数 & -- & -- & -- \\
年齢（中央値） & -- & -- & -- \\
WHO-5（開始時） & -- & -- & -- \\
推し疲れ（開始時） & -- & -- & -- \\
\bottomrule
\end{tabular}
\end{table}

\subsection{各群の前後比較（RQ2）}
\todo{要記入：表\ref{tab:prepost}に各群の開始時と6週目のWHO-5・推し疲れ，検定結果，効果量を記入}

\begin{table}[t]
\centering
\jetcaption{各群の開始時と6週目の比較（実験後に記入）}{Baseline vs.\ week 6 by group (to be completed)}
\label{tab:prepost}
\small
\begin{tabular}{@{}llccccc@{}}
\toprule
群 & 指標 & 開始時 & 6週目 & 差 & $p$（Holm） & $d_z$ \\
\midrule
A & WHO-5 & -- & -- & -- & -- & -- \\
A & 推し疲れ & -- & -- & -- & -- & -- \\
B & WHO-5 & -- & -- & -- & -- & -- \\
B & 推し疲れ & -- & -- & -- & -- & -- \\
C & WHO-5 & -- & -- & -- & -- & -- \\
C & 推し疲れ & -- & -- & -- & -- & -- \\
\bottomrule
\end{tabular}
\end{table}

\subsection{群の比較（RQ2）}
\todo{要記入：ANCOVAの結果（群Aを基準とした群B・Cの差と95\%信頼区間），傾向スコアで重み付けした群Bと群Cの比較}

\subsection{プランの実行と変化（RQ2）}
\todo{要記入：群Bのプランの実行率の分布，実行率と変化量の関係，プランごとの選ばれた回数と実行率}

\subsection{気づきと有用性（RQ3）}
\todo{要記入：気づき・使いやすさ・有用性の回答の分布，群Bと群Cの比較，自由記述の分析}

\section{考察}
\todo{実験実施後に記入：仮説H1〜H4の検証結果，予想と異なった点とその理由，正解データによる想定改善量の更新結果}
